\documentclass[11pt]{article}
\usepackage{amsmath,amssymb,amsthm,mathtools}
\usepackage[T1]{fontenc}
\usepackage[margin=1in]{geometry}
\usepackage{enumitem}
\usepackage{booktabs}
\usepackage{longtable}
\usepackage{array}
\usepackage{microtype}

\newtheorem{theorem}{Theorem}
\newtheorem{proposition}[theorem]{Proposition}
\newtheorem{remark}[theorem]{Remark}
\newtheorem{definition}[theorem]{Definition}

\title{A cleanup note on the odd-$m$ $d=5$ globalization theorem}
\author{}
\date{2026-03-14}

\begin{document}
\maketitle

\begin{abstract}
This note rewrites the current odd-$m$ $d=5$ theorem chain into a compact
manuscript-facing form. It preserves the scope distinctions emphasized by the
accepted package: the concrete bridge $(\beta,\delta)$ is \emph{componentwise}
before globalization, the exceptional landing statement is initially only a
chart/interface theorem, and the final global theorem is obtained only after
combining those inputs with the first-exit and tail-length reductions. An
appendix records the exact 26-step $m=5$ witness that is included in the bundle
as a separate finite example.
\end{abstract}

\section{Statement and provenance}

This note is a cleanup draft, not a new theorem search. It compresses the
accepted odd-$m$ $d=5$ chain into manuscript order using the files
\texttt{033}, \texttt{062}, \texttt{076}, \texttt{077}, \texttt{079},
\texttt{081}, \texttt{082}, and \texttt{083}. The theorem statement is kept in
exactly the same scope as the bundle: it is asserted for odd $m$ in the
\emph{accepted D5 regime}. We deliberately retain that phrase because the
handoff bundle states the theorem in that form and does not reduce it to one
standalone numeric cutoff.

The theorem-level object is the abstract bridge $(\beta,\rho)$, while the
concrete verifiable coordinate used for globalization is
\[
\delta = q + m\sigma,
\qquad
\sigma \equiv w + u - q - 1 \pmod m.
\]
The final claim is that the concrete bridge becomes global on the true
accessible boundary union.

\section{Imported structural inputs}

Fix the best endpoint seed
\[
L=[2,2,1],\qquad R=[1,4,4].
\]
We use the following structural inputs.

\begin{proposition}[Exact trigger family]
The target hole family on the active best-seed branch is
\[
H_{L1}=\{(q,w,u,\lambda)=(m-1,m-1,u,2):u\neq 2\}.
\]
\end{proposition}

\begin{proposition}[Universal first-exit targets]
The active branch has exactly two possible first exits to $H_{L1}$:
\[
T_{\mathrm{reg}}=(m-1,m-2,1),
\qquad
T_{\mathrm{exc}}=(m-2,m-1,1).
\]
Regular source families first exit at $T_{\mathrm{reg}}$, while the exceptional
source family first exits at $T_{\mathrm{exc}}$.
\end{proposition}

The structural role of these propositions is narrow but decisive: they pin down
the actual exceptional first exit and its branch direction before the later
chart/interface and globalization arguments are applied.

\section{Concrete boundary coordinate and componentwise bridge}

On the unique $\Theta=2$ state $(q,w,u,2)$ of a full regular chain, define
\[
\sigma \equiv w + u - q - 1 \pmod m,
\qquad
\delta := q + m\sigma.
\]

\begin{proposition}[Componentwise concrete bridge]
Under first return to the $\Theta=2$ section one has
\[
q' \equiv q+1 \pmod m,
\qquad
\sigma' \equiv \sigma + \mathbf 1_{\{q=m-1\}} \pmod m,
\]
and therefore
\[
\delta' \equiv \delta+1 \pmod{m^2}.
\]
Moreover, the current event class $\epsilon_4$ is a function of
$(\beta,\delta)$ alone. On each splice-connected accessible component, the
realized boundary labels therefore form one forward $+1$ orbit segment in
$\mathbb Z/m^2\mathbb Z$.
\end{proposition}

\begin{remark}[Scope discipline]
This proposition is intentionally \emph{componentwise}. It does not by itself
assert that $(\beta,\delta)$ is already global on the full accessible union.
That globalization is the content of the final theorem below.
\end{remark}

\section{Globalization reductions}

\begin{proposition}[Tail-length reduction]
Once the concrete bridge package $(\beta,\delta)$ is accepted, two realizations
of the same boundary label $\delta$ can differ only by remaining full-chain tail
length. Equivalently, the statement
\[
\rho = \rho(\delta)
\]
holds globally if and only if fixed realized $\delta$ has realization-independent
tail length.
\end{proposition}

\begin{proposition}[Regular continuation]
Every regular realization of every realized $\delta$ continues on the true
larger regular accessible union. Hence the full globalization theorem cannot
fail for a regular reason.
\end{proposition}

\begin{proposition}[Exceptional interface theorem]
At chart/interface level, the exceptional continuation is pinned to
\[
3m-3 \to 3m-2 \to 3m-1.
\]
The labels $3m-2$ and $3m-1$ have the distinguished regular interpretations of
terminal regular source-$4$ occurrence and regular source-$1$ start,
respectively.
\end{proposition}

\begin{remark}[Another scope boundary]
The exceptional interface proposition is a chart/interface theorem, not yet a
raw actual-lift continuation theorem. In the final proof it must be paired with
the universal first-exit theorem above.
\end{remark}

\section{Final theorem}

\begin{theorem}[Odd-$m$ D5 final globalization theorem]
For odd $m$ in the accepted D5 regime, every actual lift of the exceptional
cutoff row $\delta=3m-3$ continues through
\[
3m-2 \to 3m-1
\]
and lies in the regular continuing endpoint class there.
Consequently:
\begin{enumerate}[label=\textup{(\arabic*)},leftmargin=2.2em]
\item there is no mixed-status realized $\delta$;
\item fixed realized $\delta$ determines tail length;
\item fixed realized $\delta$ determines endpoint class and future word;
\item $\rho = \rho(\delta)$ globally;
\item raw global $(\beta,\delta)$ is exact on the true accessible boundary
union.
\end{enumerate}
\end{theorem}

\begin{proof}
By the tail-length reduction, once the concrete bridge package is in place it is
enough to show that every actual accessible realization has infinite forward
tail. The regular sector is already closed by the regular continuation theorem,
so no regular full chain is terminal.

Exceptional chains away from the cutoff continue internally by the componentwise
odometer splice law $\delta \mapsto \delta+1$. Thus only the exceptional cutoff
row $\delta=3m-3$ needs separate treatment.

For that row, the universal first-exit theorem gives a unique actual raw exit
target and branch direction for every actual lift of the exceptional family.
Determinism then forces the same immediate post-exit continuation for every such
lift. The exceptional interface theorem pins the corresponding chain-label
continuation to
\[
3m-3 \to 3m-2 \to 3m-1.
\]
Hence every actual lift of the exceptional cutoff row has an actual forward
successor and lands in the regular interface class.

Therefore every accessible full chain has a forward successor. A finite
splice-connected accessible component would then have a terminal right-end
chain, which is impossible. So every accessible component is total, and every
realized state has infinite remaining tail length. The tail-length reduction now
forces fixed realized $\delta$ to determine endpoint class and future word,
which is exactly the statement $\rho=\rho(\delta)$ globally. Raw global
$(\beta,\delta)$ exactness follows.
\end{proof}

\begin{remark}[What is and is not claimed here]
This cleanup note does not upgrade the bundle beyond its own evidence. In
particular:
\begin{enumerate}[label=\textup{(\alph*)},leftmargin=2.2em]
\item the theorem is still presented as closed \emph{inside the accepted
package};
\item the phrase \emph{accepted D5 regime} is preserved verbatim rather than
expanded to a numeric statement that the bundle does not itself isolate;
\item the small support files cover different finite ranges, and the saved
regular-union check explicitly flags $m=5$ as a tiny degenerate exception on
that particular saved regular full-chain union;
\item the explicit $m=5$ witness in Appendix~\ref{app:m5witness} is a concrete
finite example, not by itself the full odd-$m$ theorem.
\end{enumerate}
\end{remark}

\appendix

\section{Exact $m=5$ witness recorded in the bundle}\label{app:m5witness}

The bundle also contains a separate explicit finite witness
(\texttt{RoundY/d5\_m5\_kempe\_witness\_26.json}). It is described there as a
``pruned 26-move Kempe-valid witness found by greedy continuation on $m=5$.''
The JSON records
\[
\texttt{m\_verified}=(5),
\qquad
\texttt{cycle\_counts\_m5}=(1,1,1,1,1),
\qquad
\texttt{sequence\_length}=26.
\]
So this file gives an exact concrete sample instance in dimension $5$ at modulus
$m=5$.

For convenience, we record the move list exactly in the notation carried by the
JSON file.

\setlength{\LTpre}{0pt}
\setlength{\LTpost}{0pt}
\begin{longtable}{@{}>{\raggedleft\arraybackslash}p{1.2cm} >{\centering\arraybackslash}p{2cm} p{9cm}@{}}
\toprule
Step & Pair & Support \\
\midrule
\endfirsthead
\toprule
Step & Pair & Support \\
\midrule
\endhead
\bottomrule
\endfoot
1 & ${(1, 3)}$ & \texttt{("line1", 0, 2, 3)} \\
2 & ${(2, 4)}$ & \texttt{("plane", 3, null, null)} \\
3 & ${(0, 2)}$ & \texttt{("line1", 1, 4, 1)} \\
4 & ${(2, 3)}$ & \texttt{("line1", 2, 0, 2)} \\
5 & ${(1, 2)}$ & \texttt{("line2", 2, (0, 4), (0, 4))} \\
6 & ${(2, 3)}$ & \texttt{("line2", 0, (0, 4), (0, 0))} \\
7 & ${(1, 3)}$ & \texttt{("line1", 3, 2, 2)} \\
8 & ${(3, 4)}$ & \texttt{("line1", 3, 0, 4)} \\
9 & ${(0, 3)}$ & \texttt{("plane", 3, null, null)} \\
10 & ${(3, 4)}$ & \texttt{("line2", 1, (0, 2), (1, 0))} \\
11 & ${(0, 4)}$ & \texttt{("line1", 1, 1, 2)} \\
12 & ${(1, 4)}$ & \texttt{("line2", 0, (2, 3), (1, 0))} \\
13 & ${(1, 4)}$ & \texttt{("line1", 3, 4, 2)} \\
14 & ${(1, 3)}$ & \texttt{("line1", 2, 4, 4)} \\
15 & ${(1, 4)}$ & \texttt{("line1", 3, 4, 0)} \\
16 & ${(1, 3)}$ & \texttt{("line2", 0, (2, 4), (0, 1))} \\
17 & ${(0, 1)}$ & \texttt{("line1", 2, 4, 2)} \\
18 & ${(1, 4)}$ & \texttt{("line2", 0, (2, 3), (1, 2))} \\
19 & ${(1, 2)}$ & \texttt{("line2", 2, (0, 4), (4, 3))} \\
20 & ${(1, 2)}$ & \texttt{("line1", 1, 4, 3)} \\
21 & ${(1, 2)}$ & \texttt{("line1", 2, 4, 1)} \\
22 & ${(0, 1)}$ & \texttt{("line2", 2, (1, 4), (0, 4))} \\
23 & ${(1, 3)}$ & \texttt{("line1", 1, 0, 4)} \\
24 & ${(1, 2)}$ & \texttt{("line2", 1, (0, 4), (2, 4))} \\
25 & ${(0, 1)}$ & \texttt{("line2", 2, (3, 4), (4, 0))} \\
26 & ${(1, 4)}$ & \texttt{("line1", 3, 4, 1)} \\
\end{longtable}

The same JSON file also records that the same 26-step sequence is \emph{not}
a uniform odd-$m$ witness: when tested at $m=7$ it yields cycle counts
\[
(5,433,1,1,3),
\]
so the explicit 26-step witness is verified only at $m=5$.

\end{document}
